\documentclass[11pt]{article}
\usepackage{setspace}
\usepackage{geometry}
\geometry{margin=1in}
\setstretch{1.15}

\begin{document}

% --- Applicant Information ---
\noindent
\textbf{Zhongli Wang} \\
Phone: +49 160 9541 6015 \\
Email: nkwangzhongli@gmail.com \\
Website: https://sites.google.com/view/zhongli-wang \\
Bielefeld, Germany \\

\vspace{0.5cm}

% --- Recipient Information ---
\noindent
\textbf{Recruitment Committee} \\
Direction des Enqu\^etes de Conjoncture et des Analyses Micro\'economiques et Structurelles (DECAMS) \\
Service des Analyses Micro\'economiques (SAMIC) \\
Banque de France \\
75001 Paris, France \\

\vspace{0.5cm}

% --- Subject Line ---
\noindent
\textbf{Application --- Economist / Econometrician} \\

\vspace{0.3cm}

\noindent Dear Members of the Recruitment Committee,\\[0.2cm]

I am writing to apply for the Economist/Econometrician position within DGSEI--DECAMS, Service des Analyses Micro\'economiques (SAMIC). 
I am currently completing my Ph.D.\ in Economics at Bielefeld University and Universit\`a Cattolica del Sacro Cuore (Marie Sk\l odowska--Curie EPOC programme). 
My research and applied work focus on microeconometric analysis of heterogeneous agents, using firm-level data where available, and building empirical pipelines that translate granular evidence into policy-relevant conclusions. 
I am particularly motivated by SAMIC's mission to analyse household and corporate behaviour---saving and indebtedness, investment, consumption, and price-setting/inflation expectations---and to support decision-making for monetary and financial stability with rigorous, transparent evidence.

During my Ph.D.\ I developed end-to-end empirical workflows in \textbf{R, Stata, and Python} for data cleaning, record linkage across sources, feature construction, model estimation, robustness checks, and reproducible reporting. 
My dissertation studies decision-making under constraints and imperfect information, and it combines structural reasoning with empirical discipline: identifying economically meaningful mechanisms, mapping them to observables, and testing implications using micro-level data and simulation-based methods. 
This experience translates naturally to SAMIC's work, where the key challenge is to extract credible behavioural signals from granular data and to communicate them clearly to internal committees and external audiences.

I am especially interested in contributing to one of SAMIC's three research streams: \textit{corporate finance}, \textit{household finance and inequalities}, or \textit{prices/inflation expectations}. 
For corporate topics, my research in industrial organization and firm dynamics has trained me to quantify behavioural responses (e.g., investment and capacity adjustment, pricing responses to shocks, heterogeneity across firms) and to build counterfactual scenarios that are relevant for policy. 
For household topics, I have strong interest in inequality-related mechanisms and in how information frictions and beliefs affect saving, debt, and consumption choices, with an emphasis on identifying distributional effects across groups. 
For inflation expectations and price-setting, my modelling background equips me to study expectation formation and learning dynamics and to connect micro evidence to macro-relevant transmission channels.

Beyond my dissertation, I have coauthored research demonstrating my ability to work with high-dimensional micro data and to develop robust analytical systems. 
\textit{Polytope Fraud Theory} (International Review of Financial Analysis) applies supervised learning to detect structured anomalies in corporate financial statements, with careful attention to feature construction, evaluation discipline, and interpretability. 
In \textit{Agentic Workflows for Economic Research}, we develop auditable LLM-assisted pipelines that improve the efficiency and reproducibility of empirical work while maintaining human oversight---experience that is useful for large-scale literature mapping, code/documentation quality, and structured research reporting. 
Together, these projects reflect a practical strength aligned with SAMIC's needs: building reliable data products and empirical analyses that can be reviewed, updated, and communicated with confidence.

I would be excited to contribute to (i) policy notes and decision-support analysis linked to monetary policy meetings, (ii) Banque de France dissemination channels (working papers, bulletins, blog posts), and (iii) academic outputs targeting leading journals. 
I work well in collaborative environments, value careful internal review, and enjoy translating technical results into clear messages for non-specialist stakeholders. 
While my French is currently at a basic working level, I am comfortable operating professionally in English and I am actively improving my French for day-to-day work and integration in Paris.

Thank you so much for considering my application! 
I would welcome the opportunity to discuss how my microeconometric skills, programming experience, and research orientation can contribute to SAMIC's analytical and policy mission.

\vspace{6mm}

\noindent Sincerely,\\[12mm]
\noindent Zhongli Wang

\end{document}